\section{Domination patterns}
\label{sec:domination}

Domination compares a searched defender reply with an omitted reply over a
finite stopped horizon.  The comparison transfers complete strategies and
checks occupancy, window masks, and the legality frontier after every mapped
placement.  This prevents local mask similarity from hiding a different legal
continuation.

For a cell $x$, write $B_R(x)=\{z:\dist(x,z)\le R\}$ and let $\Omega(x)$ be
the 18 windows containing $x$.  For a post-opening position $P$, define its
geometric support by
\[
  \Lambda(P)=\bigcup_{s\in\Stones(P)}B_8(s).
\]
If $P$ is nonterminal, its legal placements are
$L(P)=\Lambda(P)\setminus\Stones(P)$; if $P$ is terminal, $L(P)=\varnothing$.
An empty cell is \emph{dead} when every window in $\Omega(x)$ is dead.

\subsection{Stopped outcomes and causal simulation}

The outcome records the first terminal placement, with the defender's reply
itself at time zero.

\begin{definition}[Stopped horizon outcome; companion D3 and D4]
\label{def:Dom-D3}
Let $S$ be a possibly terminal successor of a defender reply.  A horizon $n$
allows at most $n$ further placements after that reply.  Given pure
continuation strategies $\sigma_D$ and $\tau_A$, play stops at the first
terminal placement or after those $n$ placements, and
\[
  \operatorname{Out}_n(S,\sigma_D,\tau_A)
  \in\{D@t:0\le t\le n\}\cup
      \{A@t:0\le t\le n\}\cup\{?\}.
\]
Here $? $ means that no terminal placement occurred through the horizon.  For
defender comparison, set
\[
  u_D(D@t)=1,
  \qquad u_D(?)=0,
  \qquad u_D(A@t)=-1.
\]
Winning windows and times remain in the trace, although all wins by one player
have the same qualitative value.  A pure $X$-strategy assigns a legal
placement to every nonterminal history of length less than $n$ at which $X$
is to place, including both consecutive histories of an $X$ turn, and is
never queried after a terminal history.  Write $\Sigma_X(S)$ for the set of
such strategies from $S$.
\end{definition}

The transfer directions encode the defender-pruning implication.

\begin{definition}[$n$-outcome domination; companion D5 and Lemma 3]
\label{def:Dom-D5}
Let $a$ and $b$ be legal defender replies at the same nonterminal position
$P$, with stopped successors $S_a=P+a$ and $S_b=P+b$.  The reply $b$ is
\emph{$n$-outcome-dominated by $a$ for the defender}, written
$b\preceq_n a$, if there are maps
\[
  T_D:\Sigma_D(S_b)\to\Sigma_D(S_a),
  \qquad
  T_A:\Sigma_D(S_b)\times\Sigma_A(S_a)\to\Sigma_A(S_b)
\]
such that, for every discarded-branch defender strategy $\sigma_b$ and every
searched-branch attacker strategy $\tau_a$,
\[
 u_D\!\left(\operatorname{Out}_n
     (S_a,T_D(\sigma_b),\tau_a)\right)
 \ge
 u_D\!\left(\operatorname{Out}_n
     (S_b,\sigma_b,T_A(\sigma_b,\tau_a))\right).
\]
Equivalently,
\[
  \forall\sigma_b\ \exists\sigma_a\
  \forall\tau_a\ \exists\tau_b:\quad
  u_D(\operatorname{Out}_n(S_a,\sigma_a,\tau_a))
  \ge u_D(\operatorname{Out}_n(S_b,\sigma_b,\tau_b)).
\]
Consequently, an attacker win forced by the horizon after the searched reply
$a$ is also forced after the omitted reply $b$.
\end{definition}

A causal simulation is the checkable witness used by the patterns below.

\begin{definition}[Causal alternating simulation; companion D6]
\label{def:Dom-D6}
A depth-$k$ defender-favouring alternating simulation relates a searched state
$G$ to an omitted state $H$ at equal placement depth.  A defender-terminal
$G$, or an attacker-terminal $H$, is accepted immediately.  An
attacker-terminal $G$ is accepted only if $H$ is already attacker-terminal;
a defender-terminal $H$ is accepted only if $G$ is already
defender-terminal.  No move is generated after a terminal state.  Two
nonterminal states are accepted when $k=0$.

Otherwise the states have the same player to move and the same remaining
placement budget.  At a defender history, every legal move in $H$ is mapped
causally to a legal move in $G$.  At an attacker history, every legal move in
$G$ is mapped causally to a legal move in $H$.  The successor pair is related
at depth $k-1$.  The map depends only on the paired prefix and current move.
It pairs \textup{FirstStone} with \textup{FirstStone} and
\textup{SecondStone} with \textup{SecondStone}, including stored first-cell
coordinates under any coordinate involution.
\end{definition}

Such a simulation implies $n$-outcome domination by recursively transferring
the two players' strategies \cite{hexobot-companions}.  Each simulation step
must account for three channels: the mapped cells must be empty with explicit
images for branch-exclusive empties; after each mapped placement, every
labelled mask in the union of the two incident 18-window sets must be compared,
with the terminal clauses applied before phase advance; and legal moves must
be included in the actor-dependent direction.

\subsection{Dead cells and interchangeable hits}

A dead empty has no legality frontier because the dead-window witnesses on
the six directed rays already support its whole radius-eight ball.

\begin{lemma}[A dead empty has no legality frontier; companion Lemma 7]
\label{lem:Dom-L7}
Let $b$ be empty and dead in a nonterminal post-opening position $P$.  Then
\[
  B_8(b)\subseteq\Lambda(P).
\]
In particular, $b$ is legal, and placing at $b$ inserts no previously
unsupported cell into the legal store.
\end{lemma}

\begin{proof}
Let $u$ range over the six directed axial unit vectors.  The endpoint window
\[
  W_u=\{b,b+u,\ldots,b+5u\}
\]
belongs to $\Omega(b)$.  It is dead while $b$ is empty, so among
$b+u,\ldots,b+5u$ there is an old stone
$s_u=b+k_u u$ with $1\le k_u\le5$.

Take $z\in B_8(b)$.  It lies in a closed $60^\circ$ sector between adjacent
directed axes $u$ and $v$, so
\[
  z-b=\alpha u+\beta v,
  \qquad \alpha,\beta\ge0,
  \qquad \alpha+\beta=\dist(b,z)\le8.
\]
After a rotation or reflection, take $u=(1,0)$ and $v=(0,1)$, and write
$s=s_u=b+ku$ with $1\le k\le5$.  Then
\[
  \dist(z,s)=
  \max\{|\alpha-k|,\beta,|\alpha+\beta-k|\}\le8,
\]
because $\alpha$, $\beta$, $\alpha+\beta$, and $k$ all lie in $[0,8]$.
Thus $z\in B_8(s)\subseteq\Lambda(P)$.  This holds for every $z$.
\end{proof}

Deadness makes the omitted reply frontier-inert, but it does not impose the
same condition on the searched substitute.

\begin{proposition}[Frontier-certified dead-cell dismissal; companion P1]
\label{prop:Dom-P1}
Let $P$ be a nonterminal post-opening position with the defender to place.
Let $a\ne b$ be legal empty cells.  Assume every window in $\Omega(b)$ is
already dead in $P$.  Unless $P+a$ is an immediate defender win, also assume
\[
  \Lambda(P)\cup B_8(a)=\Lambda(P)\cup B_8(b).
\]
Then $b\preceq_n a$ for every finite $n$.  By \cref{lem:Dom-L7}, the displayed
condition is equivalent here to $B_8(a)\subseteq\Lambda(P)$.  Thus $b$ may be
dismissed in favor of any searched $a$ that either wins immediately or is
frontier-inert.  Deadness of $b$ does not make this separate obligation on
$a$ automatic.
\end{proposition}

\begin{proof}[Proof idea; complete proof in the companion]
Compare $P+a$ and $P+b$ by transposing $a$ and $b$ until the counterpart is
used.  The hypotheses synchronize legal support.  Every window through $b$
is permanently dead, while the remaining changed windows either agree or
change ownership in the defender-favouring direction.  The occupancy,
terminality, and frontier channels therefore give an all-depth causal
simulation; an immediate defender win after $a$ is already maximal.  See
\cite{hexobot-companions}.
\end{proof}

Two hits of one count-$4$ threat become symmetric only after every other
incident window is made inert and the successor supports agree.

\begin{proposition}[Dead-spoke interchangeable hits; companion P2]
\label{prop:Dom-P2}
Let $P$ be a nonterminal post-opening position with the defender to place.
Let $W$ be an attacker-alive window with exactly four attacker stones and
exactly two empty cells $x\ne y$.  Both $x$ and $y$ are legal.  Assume
\[
  (\Omega(x)\cup\Omega(y))\setminus\{W\}
\]
consists entirely of windows already dead in $P$, and assume
\[
  \Lambda(P)\cup B_8(x)=\Lambda(P)\cup B_8(y).
\]
Then the replies are mutually outcome-dominating at every horizon:
\[
  x\preceq_n y,
  \qquad
  y\preceq_n x
  \qquad(n<\infty).
\]
A solver may retain one canonical member of $\{x,y\}$ at that defender node.
\end{proposition}

\begin{proof}[Proof idea; complete proof in the companion]
Either defender hit kills $W$.  The dead-spoke premise then makes every window
through either special cell permanently dead.  Transposing $x$ and $y$
preserves occupancy choices, the support equality preserves legal moves in
both directions, and every mask outside the dead union is identical.  This is
an outcome-preserving bisimulation at all depths.  See
\cite{hexobot-companions}.
\end{proof}

Equality of count profiles is insufficient.  Labelled empty cells and the
window-intersection hypergraph can change minimum hitting sets or
fork-creating cells even when every pair $(\cnt_A,\cnt_D)$ agrees.  A broader
equivalence would have to preserve labelled masks, incidence, occupancy, and
radius-eight support throughout the $n$-causal cone.  The dead-spoke premise
avoids that unproved classification.

\begin{figure}[tbp]
  \centering
  \begin{tikzpicture}[scale=.52]
    \HexPatch{-1/0,0/0,1/0,2/0,3/0,4/0,5/0,6/0,
      2/-1,2/1,1/1,3/-1,4/-1,4/1,3/1,5/-1}
    \HexWindow[window one]{0}{0}{1}{0}
    \draw[black!45,line width=2.2pt]
      \HexAt{2}{-1} -- \HexAt{2}{0} -- \HexAt{2}{1};
    \draw[black!45,line width=2.2pt]
      \HexAt{1}{1} -- \HexAt{2}{0} -- \HexAt{3}{-1};
    \draw[black!45,line width=2.2pt]
      \HexAt{4}{-1} -- \HexAt{4}{0} -- \HexAt{4}{1};
    \draw[black!45,line width=2.2pt]
      \HexAt{3}{1} -- \HexAt{4}{0} -- \HexAt{5}{-1};
    \HexAttacker{0}{0}\HexAttacker{1}{0}
    \HexAttacker{3}{0}\HexAttacker{5}{0}
    \HexMark{2}{0}{dot}{}\HexMark{4}{0}{dot}{}
    \HexAttacker{2}{-1}\HexDefender{2}{1}
    \HexAttacker{1}{1}\HexDefender{3}{-1}
    \HexDefender{-1}{0}
    \HexAttacker{4}{-1}\HexDefender{4}{1}
    \HexAttacker{3}{1}\HexDefender{5}{-1}
    \HexDefender{6}{0}
  \end{tikzpicture}
  \caption{A representative dead-spoke configuration.  The highlighted
  threat has four attacker stones and the two marked empty hits.  Three spoke
  directions per hit show old attacker and defender witnesses.  The P2
  premise applies this deadness condition to all 17 other windows through
  each marked cell, not only the representative spokes shown.}
  \label{fig:dead-spoke}
\end{figure}

\subsection{Same-turn commutation}

Two placements may be sorted only when both were legal before the turn and
neither possible first placement wins.

\begin{proposition}[Nonwinning-prefix turn commutation; companion P3]
\label{prop:Dom-P3}
Let $P$ be a nonterminal \textup{FirstStone} position for either player $X$.
Let $p\ne q$ be legal already in $P$, and suppose neither $P+p$ nor $P+q$ is
terminal.  Then both traces $p;q$ and $q;p$ are legal and have the same
qualitative rule outcome after the second placement.  If nonterminal, they
have identical occupancy, labelled window masks, legal support, player to
move, and \textup{FirstStone} phase.  If terminal, both terminate on the
second placement with winner $X$ and the same absolute placement count.
\end{proposition}

\begin{proof}[Proof idea; complete proof in the companion]
Initial legality makes either second cell remain legal, and the final support
$\Lambda(P)\cup B_8(p)\cup B_8(q)$ is order-independent.  Inserting the two
owner bits into every window mask commutes, including in a window containing
both cells.  The singleton-nonterminal premises exclude a first-placement
win; any joint win is therefore completed by the second placement in both
orders.  See \cite{hexobot-companions}.
\end{proof}

The live engine states need not be field-identical.  Ordered history,
last-turn metadata, serialized order, and the retained terminal
\textup{SecondStone} first-cell payload may differ.  The proposition equates
rule outcomes, not these order-sensitive features or the two intermediate
\textup{SecondStone} threat states.

\begin{figure}[tbp]
  \centering
  \begin{tikzpicture}[
      position node/.style={draw,rounded corners=2pt,minimum width=2.4cm,
        minimum height=.8cm,fill=black!3},
      every edge/.style={draw,-{Stealth[length=5pt]}}]
    \node[position node] (P) {$P$};
    \node[position node,right=3.2cm of P] (Pp) {$P+p$};
    \node[position node,below=1.7cm of P] (Pq) {$P+q$};
    \node[position node,right=3.2cm of Pq] (Ppq) {$P+\{p,q\}$};
    \path (P) edge node[above] {$p$} (Pp)
              edge node[left] {$q$} (Pq)
          (Pp) edge node[right] {$q$} (Ppq)
          (Pq) edge node[below] {$p$} (Ppq);
  \end{tikzpicture}
  \caption{The commuting square for two same-turn placements.  Both paths
  reach the same completed-turn rule state under the premises of P3.}
  \label{fig:commutation}
\end{figure}

The analytic proofs have randomized local checks, not exhaustive machine
proofs.  MV-P3 tested 2,997 random legal same-mover pairs under the
singleton-nonwin filter and found 0 commutation mismatches.  MV-L7 tested 400
adversarial dead-cell configurations and found 0 legality-frontier violations;
the worst minimum distance from a cell of the radius-eight ball to an old
stone was $6$, leaving margin $2$.  Full exhaustive quotient enumeration
remains open.

These results do not claim that an arbitrary frontier-extending reply can
substitute for a dead cell, that equal window counts make hitting cells
interchangeable, or that stones of one colour are globally monotone-helpful
when legal frontiers differ; no transferred strategy continues after a
terminal placement, and the prospective exhaustive enumerations are
specifications rather than executed evidence.
